\subsection{Kesimpulan}

Berdasarkan pembahasan pada bab-bab sebelumnya, dapat disimpulkan bahwa masyarakat perkotaan memiliki kebutuhan terhadap layanan kebersihan rumah tangga yang lebih praktis, terpercaya, transparan, dan dapat dipantau secara real-time. Model layanan konvensional masih memiliki beberapa kendala, terutama dalam hal verifikasi kredibilitas petugas, ketidakjelasan harga, dan ketidakpastian waktu kedatangan.

CleanConnect dirancang sebagai solusi teknologi berbasis aplikasi seluler yang menghubungkan pengguna dengan penyedia jasa kebersihan rumah tangga melalui fitur booking, tracking petugas berbasis lokasi/GPS, rating dan review, estimasi biaya otomatis, serta jadwal rutin mingguan atau bulanan. Fitur-fitur tersebut diharapkan dapat meningkatkan kepercayaan pengguna, memperjelas struktur biaya, dan membuat proses pemesanan layanan menjadi lebih efisien.

\subsection{Saran}

Pengembangan CleanConnect selanjutnya disarankan untuk dilanjutkan ke tahap perancangan prototipe antarmuka dan pengujian langsung kepada calon pengguna. Pengujian tersebut penting untuk memastikan bahwa alur pemesanan, tampilan estimasi biaya, fitur pelacakan petugas, serta sistem rating dan review benar-benar mudah dipahami dan sesuai dengan kebutuhan masyarakat perkotaan.

Selain itu, penelitian lanjutan perlu melibatkan data pengguna yang lebih konkret, seperti hasil wawancara, survei, atau observasi lapangan, agar rancangan sistem dapat disesuaikan dengan kondisi nyata pengguna dan penyedia jasa kebersihan.
